\section{연구 결과}

연구 문제의 순서에 맞추어 결과를 제시하세요. 과학교육 연구에서는 표와 그림이 단순 요약을 넘어 해석의 근거가 되도록, 본문에서 주요 패턴과 의미를 함께 설명하는 것이 좋습니다.

표와 그림은 학술지 규정에 맞추어 제목, 주석, 단위, 약어 설명을 일관되게 작성하세요. 본문에서 언급하지 않은 표와 그림은 넣지 않는 것이 좋습니다.

한국과학교육학회지는 최종 조판에서 2단 편집을 사용하므로, 표나 그림을 만들 때 처음부터 폭을 1.0 textwidth 또는 0.5 textwidth 중 하나로 정해 두는 것이 좋습니다. 두 단을 모두 사용하는 큰 그림, 복잡한 표, 여러 패널이 있는 결과는 1.0 textwidth를 기준으로 만들고, 한 단 안에서 충분히 읽히는 단순 그림이나 작은 표는 0.5 textwidth를 기준으로 만듭니다. 이렇게 정해 두면 최종 조판 단계에서 그림 글자 크기, 축 라벨, 범례, 표 간격이 갑자기 작아지는 문제를 줄일 수 있습니다.

\begin{ManuscriptWideTable}
\caption{한국과학교육학회지 원고용 표 예시}
\label{tab:example}
\centering
\begin{tabular}{ll}
\toprule
구분 & 내용 \\
\midrule
예시 & 연구 문제별 결과 또는 분석 범주를 입력하세요. \\
\bottomrule
\end{tabular}
\end{ManuscriptWideTable}
